\section{The finite radial decomposition without the larger normal form}
\label{rd-section}

We now construct the algebra homomorphism required in
\eqref{rp-required-normal-form}.  Only the terminal
operator proved in Section~\ref{rp-section} and the
compact branching law of Theorem~\ref{ib-compact-branching}
are needed.  In particular, the construction does not
require the additional raising and lowering operators
in the larger source normal form
\cite{alpo-uniform-source}.

\subsection{The function space and the actual matrix field}

Retain the Hermitian space $V$ of signature $(k,1)$,
the vectors $z,z'$ and the positive constants $T,T'$
from Section~\ref{rp-section}.  Let
\[
 K_z=\{a\in U(V,h):az=z\},\qquad
 K_{z'}=\{b\in U(V,h):bz'=z'\}.
\]
Let $\pi:K_z\to U(W')$ and
$\pi':K_{z'}\to U(W)$ be nonzero irreducible,
finite-dimensional continuous unitary representations.
Let $r\geq0$ and set
\[
 \Omega=(\mathbb C^\times)^r,\qquad
 D=\{g\in U(V,h):G(g)>0\}.
\]
For $r=0$, $\Omega$ is a one-point space.
Define $\mathscr R$ to be the space of smooth functions
\[
 \Phi:\Omega\times D\longrightarrow\operatorname{Hom}(W,W')
\]
with
\begin{equation}\label{rd-biequivariance}
 \Phi(c,agb)=\pi(a)\Phi(c,g)\pi'(b)
              \quad(a\in K_z,\ b\in K_{z'}).
\end{equation}
Extend these functions by zero outside $\Omega\times D$.
Smoothness is required on this open submanifold only,
not across its boundary.  This is the source's radial
function space after identifying the terminal group
with $U(V,h)$.  Its formulation using local smooth
extensions to matrix space is equivalent: local
coordinates for an embedded Lie group extend a smooth
function by making it constant in the normal
coordinates.

For $g\in D$, write $x=z'$, $y=g^{-1}z$, and let
\[
                       E_g=\langle y,x\rangle^\perp.
\]
This is a positive space of dimension $k-1$.
Using any orthonormal basis of $E_g$, define the
matrix $C_g(Y)$ on $W$ by
Corollary~\ref{rp-finite-terminal}.

\begin{lemma}\label{rd-smooth-matrix}
The matrix $C_g(Y)$ is independent of the chosen
orthonormal basis of $E_g$.  It is smooth in $g$,
polynomial in $Y$, and obeys
\begin{equation}\label{rd-matrix-equivariance}
 C_{agb}(Y)=\pi'(b)^{-1}C_g(Y)\pi'(b).
\end{equation}
Consequently
\begin{equation}\label{rd-actual-X}
             (X(Y)\Phi)(c,g)=\Phi(c,g)C_g(Y)
\end{equation}
defines an endomorphism of $\mathscr R$.  This is the
actual differential operator
$G^{-1}\mathcal T_{22}(Y)$, and it does not enlarge
the support of a function.
\end{lemma}

\begin{proof}
A change of orthonormal basis in $E_g$ conjugates its
general-linear generators exactly as in
Lemma~\ref{rp-capelli}.  Full determinant covariance
makes $C_g(Y)$ unchanged.

To check smoothness without choosing a global frame,
put
\[
 v_g=y+\frac{h(y,x)}{T'}x.
\]
Then $v_g\in x^\perp$ and direct expansion gives
$h(v_g,v_g)=G(g)/T'>0$.  The space $E_g$ is the
orthogonal complement of $v_g$ in the fixed positive
space $x^\perp$.  Its orthogonal projection is
\[
       v\longmapsto v-\frac{h(v,v_g)}{h(v_g,v_g)}v_g
                                      \quad(v\in x^\perp),
\]
which varies smoothly in $g$.  Near a fixed point,
project a fixed basis of $E_g$ by these projections
and apply positive-definite Gram--Schmidt.  The
resulting local orthonormal frame is smooth, since
all denominators and the required squared norms
are positive near that point.  The definition of
$C_g(Y)$ is polynomial in the associated rank-one
matrices and their images under $d\pi'_{\mathbb C}$,
so it is smooth in that neighborhood.  Frame
independence makes these descriptions agree.

For $g'=agb$ one has $y'=b^{-1}y$, while
$b^{-1}x=x$.  Thus $E_{g'}=b^{-1}E_g$.
Conjugating the rank-one matrices by $b^{-1}$
conjugates their representation operators by
$\pi'(b)^{-1}$.  Ordered products give
\eqref{rd-matrix-equivariance}.
Combining it with \eqref{rd-biequivariance} proves
that \eqref{rd-actual-X} again satisfies
\eqref{rd-biequivariance}.  Joint smoothness and
the zero extension are preserved.
Corollary~\ref{rp-finite-terminal}, applied to
$g\mapsto\Phi(c,g)$ for fixed $c$, identifies it
with the actual differential operator.  The
displayed multiplication formula proves the
support assertion.
\end{proof}

\subsection{Matching the source parameters with highest weights}

Let the spectral parameters $u,u'$ satisfy
\eqref{rp-interlacing}, and suppose the compact types
are specified by the source's Capelli equations:
on the representation with parameters $u$ the
full $k$-dimensional determinant is
\begin{equation}\label{rd-source-Capelli}
       Q_k(w)=\prod_{i=1}^k
           \left(w-\frac{k-1}{2}+u_i\right)\operatorname{Id},
\end{equation}
and the analogous equation holds for $u'$.
The common integral coset condition is as in
Section~\ref{rp-section}.

\begin{lemma}\label{rd-highest-parameters}
With highest weights indexed from $1$ to $k$, the
representations $\pi,\pi'$ have respectively
\begin{equation}\label{rd-lambda-u}
 \lambda_i=u_i+i-1-\frac{k-1}{2},
 \qquad
 \lambda'_i=u'_i+i-1-\frac{k-1}{2}.
\end{equation}
In particular these numbers are integers.
The common irreducible constituents on restriction
to a common positive $(k-1)$-dimensional subspace
are indexed exactly by the tuples $\rho\in\mathcal S$
from Section~\ref{rp-section}, under
\begin{equation}\label{rd-mu-rho}
                   \rho_i=\mu_i-i+1+\frac{k-2}{2}.
\end{equation}
On such a constituent the operator $C_g(Y)$ acts by
$p_\rho(Y)$.
\end{lemma}

\begin{proof}
The highest-weight classification in
Lemma~\ref{cb-highest-classification} gives integral
dominant weights $\lambda,\lambda'$.
Corollary~\ref{rp-central-character} gives the
full Capelli eigenvalue
$\prod_{i=1}^k(w-i+1+\lambda_i)$.
Its constants $\lambda_i-i+1$ are strictly
decreasing.  The $u_i$ are also strictly decreasing
by \eqref{rp-interlacing}.  Comparison with
\eqref{rd-source-Capelli} and equality of ordered
real roots prove \eqref{rd-lambda-u}.

By Theorem~\ref{ib-compact-branching}, a highest
weight $\mu$ occurs on restriction from the left
type precisely when
$\lambda_i\geq\mu_i\geq\lambda_{i+1}$, and similarly
for the right type.  Applying \eqref{rd-lambda-u}
and \eqref{rd-mu-rho}, the simultaneous conditions
become
\[
 \max(u_{i+1},u'_{i+1})+\tfrac12
          \leq \rho_i
          \leq \min(u_i,u'_i)-\tfrac12 .
\]
Moreover
$\rho_i-\max(u_{i+1},u'_{i+1})-1/2
=\mu_i-\max(\lambda_{i+1},\lambda'_{i+1})$
is an integer.  These are exactly the finite sets
defining $\mathcal S$.
The inequalities automatically make $\mu$ dominant.
Finally the $(k-1)$-dimensional Capelli eigenvalue
at $w=Y+(k-2)/2$ is
\[
 \prod_{i=1}^{k-1}(Y+(k-2)/2-i+1+\mu_i)
                                      =p_\rho(Y).
\]
For $k=1$ the common subgroup is trivial, each
original representation is one-dimensional,
$\mathcal S$ has the empty tuple, and this last
determinant is $1$.
\end{proof}

\begin{lemma}\label{rd-common-types}
For every $\Phi\in\mathscr R$ and $g\in D$,
the map $\Phi(c,g)$ vanishes on all right
constituents whose type does not also occur on
the left.  Thus its nonzero components have
precisely the possible Capelli polynomials
$p_\rho(Y)$, $\rho\in\mathcal S$.
\end{lemma}

\begin{proof}
Let $H_g$ be the subgroup that acts arbitrarily
on $E_g$ and as the identity on $\langle y,x\rangle$.
It is isomorphic to $U(k-1)$.
For $h\in H_g$, one has $h\in K_{z'}$ and
$ghg^{-1}\in K_z$.  Evaluating
\eqref{rd-biequivariance} in these two ways gives
\[
             \Phi(c,g)\pi'(h)
                 =\pi(ghg^{-1})\Phi(c,g).
\]
Both restrictions are orthogonal sums of the
irreducibles in the compact branching theorem.
An intertwiner between two nonisomorphic
irreducibles is zero: its kernel and image
are invariant, so a nonzero such map would
be an isomorphism.  Projection onto the
orthogonal left summands therefore shows that
every noncommon right summand is annihilated.
Lemma~\ref{rd-highest-parameters} identifies the
remaining labels and eigenvalues.
\end{proof}

\subsection{Construction of the projections}

\begin{theorem}\label{rd-constructed-kappa}
There is a unique unital algebra homomorphism
\[
       \kappa:\mathbb C^{\mathcal S}
                    \longrightarrow
                         \operatorname{End}_{\mathbb C}(\mathscr R)
\]
such that
\begin{equation}\label{rd-kappa-identity}
       \kappa((p_\rho(Y))_\rho)
                =X(Y)=G^{-1}\mathcal T_{22}(Y)
                            \qquad(Y\in\mathbb C).
\end{equation}
It can be chosen as a polynomial in one actual
terminal operator.  For
$P_\rho=\kappa(1_{\{\rho\}})$, every radial
function has the decomposition
\begin{equation}\label{rd-decomposition}
 \Phi=\sum_{\rho\in\mathcal S}P_\rho\Phi,\qquad
 \mathcal T_{22}(Y)P_\rho\Phi
                        =G\,p_\rho(Y)P_\rho\Phi.
\end{equation}
At every point the summands are orthogonal in
the Hilbert--Schmidt form.  In particular
\begin{equation}\label{rd-pointwise-contraction}
 \sum_\rho\|(P_\rho\Phi)(c,g)\|_{\rm HS}^2
                      =\|\Phi(c,g)\|_{\rm HS}^2 .
\end{equation}
All projections preserve the stated smooth
equivariance and do not enlarge supports.
\end{theorem}

\begin{proof}
Choose a real $Y_0$ for which
$\xi_\rho=p_\rho(Y_0)$ are distinct, as in
Proposition~\ref{rp-idempotent-positive}.
Let
\[
 q_0(Z)=\prod_{\rho\in\mathcal S}(Z-\xi_\rho),
 \qquad
 \ell_\rho(Z)=
       \prod_{\sigma\ne\rho}
           \frac{Z-\xi_\sigma}{\xi_\rho-\xi_\sigma}.
\]
We first prove $q_0(X(Y_0))=0$ on $\mathscr R$.
At a fixed $(c,g)$, induction on the degree shows
that a polynomial $a$ obeys
\[
 (a(X(Y_0))\Phi)(c,g)
                         =\Phi(c,g)a(C_g(Y_0)).
\]
This uses \eqref{rd-actual-X} for every function
to which the next operator is applied.
Lemma~\ref{rd-common-types} shows that
$\Phi(c,g)q_0(C_g(Y_0))=0$: on common types
the factor is zero, while on other types
$\Phi(c,g)$ itself is zero.
Outside the prescribed domain all functions
are zero.  Pointwise equality proves
$q_0(X(Y_0))=0$.

In the quotient algebra $\mathbb C[Z]/(q_0)$,
the $\ell_\rho$ are mutually annihilating
idempotents and their sum is $1$.
Indeed each claimed polynomial identity
vanishes at every distinct $\xi_\rho$ and
is therefore divisible by $q_0$.
Set
\[
     P_\rho=\ell_\rho(X(Y_0)),\qquad
     \kappa(f)=\sum_\rho f(\rho)P_\rho.
\]
The same identities hold as operators, proving
that $\kappa$ is a unital algebra homomorphism.
For a right constituent of common type $\sigma$,
its evaluation satisfies
\[
            (P_\rho\Phi)(c,g)
              =\Phi(c,g)\ell_\rho(C_g(Y_0)),
 \qquad
 \ell_\rho(C_g(Y_0))|_{W_\sigma}
                          =\delta_{\rho,\sigma}\operatorname{Id}.
\]
Noncommon constituents are killed by $\Phi(c,g)$.
Consequently both sides of
\eqref{rd-kappa-identity}, evaluated at $(c,g)$,
act as $p_\sigma(Y)$ on every surviving
constituent, and agree.  This proves the identity
for all complex $Y$, without an analytic
continuation argument.

The idempotent identities now imply
\eqref{rd-decomposition}.  Since $C_g(Y_0)$ is
self-adjoint, the proof of
Proposition~\ref{rp-idempotent-positive} gives
pointwise orthogonality and
\eqref{rd-pointwise-contraction}.
Lemma~\ref{rd-smooth-matrix} shows that every
polynomial in $X(Y_0)$ acts within $\mathscr R$.
Its pointwise matrix formula also preserves
the support of the original function.

Finally any unital algebra homomorphism satisfying
\eqref{rd-kappa-identity} must send $1_{\{\rho\}}$
to $\ell_\rho(X(Y_0))$, since
$1_{\{\rho\}}=\ell_\rho((\xi_\sigma)_\sigma)$
in the finite function algebra.  This proves
uniqueness.  In particular its action agrees
with the action of the source's larger normal
form whenever that form is supplied.
\end{proof}

No estimate on the separation between the
$\xi_\rho$ is needed.  They are used to construct
the projections, but the resulting estimate
\eqref{rd-pointwise-contraction} has constant one.
In particular there is no hidden dependence on
the size of the type set or on the interpolation
denominators in this estimate.

\begin{corollary}\label{rd-integrated-projections}
Let $\nu$ be any positive measure on
$\Omega\times D$.  Whenever the functions in
question are square integrable, the projections
$P_\rho$ are mutually orthogonal contractions
for the Hilbert--Schmidt $L^2(\nu)$ norm, and
\[
 \|\Phi\|_{L^2(\nu)}^2
                  =\sum_\rho\|P_\rho\Phi\|_{L^2(\nu)}^2.
\]
For $f\in\mathbb C^{\mathcal S}$,
\[
 \|\kappa(f)\Phi\|_{L^2(\nu)}
              \leq \max_\rho|f(\rho)|\,\|\Phi\|_{L^2(\nu)}.
\]
The same assertions hold after inserting any
nonnegative measurable scalar weight into the
measure.  Each $P_\rho$ commutes with scalar
multiplication and with differentiation in
the variables $c\in\Omega$.
\end{corollary}

\begin{proof}
Integrate \eqref{rd-pointwise-contraction};
the sum is finite.  More generally, its
orthogonal components give
$\|\kappa(f)\Phi(c,g)\|_{\rm HS}^2
=\sum_\rho|f(\rho)|^2
\|(P_\rho\Phi)(c,g)\|_{\rm HS}^2$.
Integration proves the bound.
The projections are pointwise right
multiplications by matrices depending only
on $g$, which proves the last assertions
whenever both sides are defined on the
smooth radial space.
\end{proof}

\begin{corollary}\label{rd-actual-positive-density}
Under the actual compact-type and spectral-margin
hypotheses, the terminal density
\eqref{rp-terminal-density} is nonnegative
without an additional normal-form premise.
For a fixed $c$, it is integrable against Haar
measure on $U(V,h)$ if all its finitely many
functions $\Psi_e(c,\cdot)$ have compact
support in $D$.
\end{corollary}

\begin{proof}
Theorem~\ref{rd-constructed-kappa} supplies
exactly \eqref{rp-required-normal-form}, so
Corollary~\ref{rp-terminal-density-positive}
applies.  On the finite union of the compact
supports, $G$ has a positive minimum and
$|\zeta|^2/G$ is bounded.  The functions,
their projections, and all displayed
coefficients are continuous there.
The density is therefore bounded and
compactly supported.  Haar measure is
finite on compact sets, proving
integrability.
\end{proof}

The construction proves the entire finite
spectral decomposition needed for terminal
positivity, not just a conditional implication
from the larger normal form.  It also proves
an exact contraction estimate for arbitrary
positive measures.  The remaining analytic
task is to identify the original integrated
differential forms with these densities on
the source's generated function classes,
including their noncompact members and
successive radial integrations.  Compact
support at one fixed value of $c$ is not
asserted for that whole class.


